% SPDX-License-Identifier: GPL-3.0-only
% Licensed under GNU GPL version 3; https://www.gnu.org/licenses/gpl-3.0.html . See LICENSE at repository root.
\section{Worked examples and semantic boundaries}\label{sec:examples}

\subsection{An assignment changes the next-state observation}
Use the integer model, a variable $x$, and the constant $1$. Let $U$ be any
rooted state expression, $U'=\upd{x:=1}{U}$, and $P=(x=1)$. The goal is
\[
 \seq{}{\ann{\epsilon\jlabel U}{[x:=1]\Next P}}.
\]
There is no initial assumption on $x$. The six-node acyclic proof proceeds
as follows; the tag on the box is suppressed in the display.
\begin{enumerate}
\item Assignment-right gives
$\seq{}{\ann{(\epsilon\jlabel U)\jlabel U'}{\Next P}}$.
\item An explicit unit conversion presents
$\seq{}{\ann{U\jlabel U'}{\Suf{\bot}{P}}}$.
The generator does not silently normalize away the empty prefix.
\item Suffix-right produces two premises, in order:
\[
 \seq{}{\ann{U'}{P},\ann{U'}{\bot}},\qquad
 \seq{}{\ann{U'}{P},\ann{U'}{\Suf{\bot}{P}}}.
\]
\item Close the first premise with a supplied pure proof: at $U'$, $x=1$.
On the second premise, weaken away the recursive suffix occurrence, giving
$\seq{}{\ann{U'}{P}}$, and close that atom/equality leaf with the same
state-update calculation.
\end{enumerate}
The suffix step contributes two nodes, so the graph contains six nodes in
total. No infinite branch exists, and the empty trace witness suffices. The
exported theorem follows from whole-graph acceptance through
Theorem~\ref{thm:certificate}.
\sources{\rocq{Benchmarks/AssignmentChange.v}{assignment_changes_next_value},
  \rocq{Benchmarks/AssignmentChange.v}{first_suffix_edge_has_no_progress},
  \rocq{Benchmarks/AssignmentChange.v}{continuing_suffix_edge_progresses}}

This example also explains why $[x:=1](x=1)$ expresses a different claim.
The atom in that formula still reads the original first state. Moving to
the next suffix is what observes the state written by the assignment.

\subsection{A genuine cyclic certificate preserves the first state}
Let $P=(x\ge0)$, $\alpha=(x:=x+1)$, and let $\ell_0$ be a singleton state
label, $k$ a nonempty trace name, and $V$ a final state expression. The
benchmark's root has shape
\[
 \seq{\ann{\ell_0}{P}}
   {\ann{(\ell_0\jlabel k)\jlabel V}{[\alpha^*]\Fst{P}}^{0}}.
\]
The initial state expression in $\ell_0$ includes a parallel update to
another variable. This does not affect the argument: the same initial
label is shared by both sides.

Star-right gives a finite base branch and a recurrent branch. On the base
branch, expand the true test, reassociate the label explicitly, apply
first-right, and close by identity. On the recurrent branch, assignment-right
appends the updated final state. A label rotation exposes a trace
substitution that absorbs the old final state into $k$, followed by a world
substitution for the new final state. The resulting premise is the original
root, and an exact backlink connects it to node zero.
\sources{\rocq{Benchmarks/Cyclic.v}{star_prefix_invariant}}

The live cycle and one accepted witness are
\[
 0\longrightarrow2\mathrel{\mathop{\longrightarrow}\limits^{\text{assignment}}}7
 \longrightarrow8\longrightarrow9\longrightarrow10\longrightarrow0,
\]
\begin{center}
\begin{tabular}{@{}lrrrrrr@{}}
\toprule
Node & 0 & 2 & 7 & 8 & 9 & 10\\
\midrule
Selected tag & 0 & 0 & 0 & 0 & 0 & 0\\
Potential & 2 & 1 & 6 & 5 & 4 & 3\\
\bottomrule
\end{tabular}
\end{center}
The assignment edge strictly decreases the semantic rank, allowing the
potential to reset from $1$ to $6$. Every other cycle edge decreases the
potential, including the backlink from potential $3$ to $2$. The finite base
branch needs no followed tag. The graph has eleven nodes. A separate
twelve-node benchmark inserts a checked tag-renaming backlink.
\sources{\rocq{Benchmarks/Cyclic.v}{witness},
  \rocq{Benchmarks/Cyclic.v}{graph_size},
  \rocq{Benchmarks/CyclicTags.v}{explicit_tag_renaming_cycle}}

The cyclic theorem is a property of the \emph{first} observed state. It does
not claim that the increment program preserves an arbitrary property of
every visited state; its use of $\Fst{P}$ is essential to the displayed claim.

\subsection{A guarded loop is safe at every visited state}
In the integer model let
\[
 \alpha=(0<x)?;(x:=x-1),\qquad I(s)\Longleftrightarrow s(x)\ge0.
\]
The guard/test/assignment execution has trace $[s,s,s']$, where
$s'(x)=s(x)-1$. If $I(s)$ and the test succeeds, then $I$ holds at all three
states. Induction over arbitrary finite powers lifts this property to every
run of $\alpha^*$, including the final stutter. The pure-always theorem then
gives
\[
 \seq{\ann{U}{x\ge0}}{\ann{U}{[\alpha^*]\Alw{(x\ge0)}}}.
\]
This benchmark is a direct semantic invariant proof for all finite iteration
counts, rather than a cyclic certificate.
\sources{\rocq{Benchmarks/CounterLoop.v}{step_safe},
  \rocq{Benchmarks/CounterLoop.v}{every_finite_stop},
  \rocq{Benchmarks/CounterLoop.v}{guarded_counter_all_visited_nonnegative}}

There is no claim of reaching zero: even from $x=7$, the zero-iteration trace
$[s,s]$ is a valid star run. Safety at all visited states and termination at
an exit state are different properties under these finite-star semantics.
\sources{\rocq{Benchmarks/CounterLoop.v}{may_stop_before_exit}}

\subsection{A user-defined heap theory}\label{sec:heap-example}
Take the three carriers from Section~\ref{sec:signature}, and define
\begin{align*}
 \mathsf{read}(h,i)&=h(i),\\
 \mathsf{store}(h,i,v)(j)&=
   \begin{cases}v,&j=i,\\h(j),&j\ne i.\end{cases}
\end{align*}
The supplied model proves
\[
 \mathsf{read}(\mathsf{store}(h,i,v),i)=v,
 \qquad j\ne i\ \Longrightarrow\
 \mathsf{read}(\mathsf{store}(h,i,v),j)=\mathsf{read}(h,j).
\]
No equality decision on the heap carrier $\Z\to\Z$ is required.
\sources{\rocq{Benchmarks/UserStructures.v}{read_store_same},
  \rocq{Benchmarks/UserStructures.v}{read_store_other}}

Let $h:\mathsf{heap}$ and $i:\mathsf{int}$ be program variables. The
certified assignment example proves
\[
 \seq{}{\ann{\epsilon\jlabel U}
   {[h:=\mathsf{store}(h,i,1)]\Next(\mathsf{read}(h,i)=1)}}.
\]
Its graph has the same assignment/unit/suffix/weakening structure as the
scalar example. At a pure leaf, the user supplies the proved read-after-store
law. The declaration of \code{Store} alone would not discharge that leaf.
The same model separately certifies quantifier reflexivity at heap sort.
\sources{\rocq{Benchmarks/UserStructures.v}{heap_assignment_reads_written_value},
  \rocq{Benchmarks/UserStructures.v}{quantified_heap_reflexivity}}

\subsection{Small calculations that expose the boundaries}
The following checks are useful when reading or extending the theory.
\begin{enumerate}
\item On $[s_1,s_1,s_0]$, if $P$ means $x=0$, $s_1(x)=1$, and $s_0(x)=0$,
$\Suf{\top}{P}$ holds but $\Next P$ does not. The suffix operator may postpone
its choice; next cannot.
\item $\Fst{\Next\top}$ is false on every nonempty trace, because the first
operator evaluates its body on a singleton. In particular, it is false on
$[s,s]$ even though $\Next\top$ is true on that two-state trace.
\item For three accumulated traces of lengths $1,2,3$, the history measure
is $6$, while the number of history components is $3$. Counting executions
or transitions would give a different rank.
\item Sequential and parallel assignments can have different results; the
snapshot swap in Section~\ref{sec:updates} is checked in the benchmark.
\item A variable fresh in formulas can still be stale in a label. The
quantifier benchmark gives a valid proposed premise and a false conclusion
when that forbidden eigenvariable choice is attempted; the request rejects it.
\end{enumerate}
\sources{\rocq{Benchmarks/Temporal.v}{suffix_postponed},
  \rocq{Benchmarks/Temporal.v}{next_cannot_postpone},
  \rocq{Benchmarks/Temporal.v}{compound_first_state},
  \rocq{Benchmarks/Temporal.v}{sum_of_component_lengths},
  \rocq{Benchmarks/UpdateAlgebra.v}{parallel_swap},
  \rocq{Benchmarks/Quantifiers.v}{stale_label_eigenvariable_rejected}}
